\begin{textsample}{NEG Dim 6 – vem\_ed\_15 – Score 1.00 – vem\_ed\_15/t065.txt}  \label{ex:f6_neg_140}
PERSPECTIVAINTERNACIONAL A parte da manhã (9h30 às 12h) foi dedicada aos relatos de parceiros internacionais das Fatecs na realização de PCIs.
Da China, as professoras Sally Xu e Hang Yin (Tianjin Normal University) compartilharam a experiência de um PCI desenvolvido há quatro anos com Fatecs, com uma participação de aproximadamente 1.000 alunos brasileiros e chineses por ano.
Voltado às competências interculturais e ao aprendizado do inglês por parte de falantes não-nativos, envolveu nove unidades em 2022: Campinas, Catanduva, Franca, Ferraz de Vasconcelos, Itaquaquecetuba, Mogi das Cruzes, Piracicaba, São José do Rio Preto e Sebrae.
Adya Sharma, diretora do Symbiosis Centre for Management Studies (Symbiosis International, Deemed University)/Índia, apresentou o país, a instituição e os PCIs / COIL desenvolvidos com Fatecs.
Os professores Sonica Rautela, Nehajoan Panackal e Ashutosh Mathur comentaram o desenvolvimento de um estudo comparativo na área de marketing.
Estudantes da IES indiana também gravaram depoimentos, celebrando a aprendizagem sobre novas culturas, a interação e a ampliação de horizontes.
Divinia Jithoo, especialista em educação internacional na Durban University of Technology (DUT / África do Sul), comentou sobre a forte parceria entre DUT e as Fatecs.
Como são duas instituições no eixo de cooperação Sul-Sul, os estudantes enfrentam desafios similares e, devido às limitações financeiras, a realização de intercâmbios virtuais torna as ações pedagógicas mais viáveis.
Lindewa Mkhize, coordenadora de educação internacional da DUT, participou ao vivo da sessão de perguntas e destacou a colaboração entre coordenadores, professores e alunos.
Hope Windle, diretora do SUNY COIL Center (EUA), contou como a iniciativa COIL se expandiu pelo mundo, incluindo as Fatecs do Centro Paula Souza.
Trata-se de um relacionamento de longa data, que iniciou com um projeto realizado em 2013 entre SUNY Ulster e Fatec Americana.
Entre as principais lições aprendidas estão: ouvir os parceiros, ser flexível, definir os papéis de cada um e refletir sobre a experiência.
Javier Guerrero, coordenador de desenvolvimento de professores da Uniminuto (Colômbia), comentou sobre os projetos realizados desde 2021 com Fatecs, envolvendo áreas de administração e licenciaturas em inglês e espanhol.
Ressaltou a transferência de conhecimento pela equipe dos PCIs na realização de intercâmbios virtuais e capacitação de professores da Uniminuto.
A professora Bianca Cely mencionou a importância da realização de eventos acadêmicos em conjunto como o 1ºColóquio de Literatura Brasileira, realizado em 2021com a participação de Regiane Camargo (equipe dos PCIs / Cesu) e Wilton Garcia (Fatec Itaquaquecetuba).

% matched lemmas: 
\end{textsample}
